\begin{textsample}{POS Dim 1 – go – Score 42.00 – go/go\_inf\_20.txt}  \label{ex:f1_pos_104}
3.4.
Escuta, Fala, Pensamento e Imaginação Toda criança, desde bebê, tem o direito de conhecer, experienciar, compreender e se apropriar da sua língua materna, tanto no \textbf{que} se refere à linguagem oral, falada ou por sinais ( Língua Brasileira de Sinais - Libras ), quanto no \textbf{que} \textbf{diz} respeito à linguagem escrita, convencional ou tátil ( Braile ), ou ainda, por outros sistemas simbólicos \textbf{que} têm como suporte a tecnologia assistiva.
A língua como instrumento de comunicação se caracteriza como atividade humana carregada de sentidos e significados compartilhados por determinado grupo social, \textbf{que} possibilita a constituição do \textbf{sujeito} e da sua consciência, bem como o desenvolvimento do pensamento e da imaginação.
A língua é \textbf{um} sistema simbólico \textbf{que} tem como função mediar a relação dos \textbf{sujeitos} com o mundo e com eles mesmos.
Ela se apresenta por meio de signos, de regras, de convenções e de relações \textbf{que} constituem as funções psíquicas humanas, como a \textbf{percepção}, a memória, a atenção, a ação voluntária, a comunicação, o conhecimento, o \textbf{raciocínio}, a reflexão, a autorregulação do comportamento, a consciência e a afetividade.
Portanto, sua apropriação é \textbf{um} processo complexo \textbf{que} \textbf{necessita} de mediação.
A compreensão e o uso da língua pela criança em diferentes situações comunicativas, entendidas como espaços de negociação de significados \textbf{que} levem à compreensão do \textbf{que} se \textbf{diz}, possibilita \textbf{que} ela desenvolva a sua capacidade de simbolizar, diversificando as suas maneiras de entender e interpretar a realidade.
A BNCC ( BRASIL, 2017 ) apresenta este campo de experiências abordando conceitos \textbf{centrais} \textbf{que} permitem apreender o processo de apropriação da língua materna pela criança, bem como suas formas de expressar leituras, compreensões, \textbf{percepções} e sentimentos sobre o mundo, sobre os outros e sobre si mesma, conforme segue : É \textbf{inerente} a esse campo de experiências a discussão sobre \textbf{um} dos grandes desafios da Educação Infantil, refletir e entender qual é o papel desta etapa da Educação Básica no processo de compreensão, uso e apropriação da língua materna pela criança.
Dessa forma, este campo de experiência tem por objetivo abordar a articulação entre a escuta, a fala[^19 ] e o pensamento, o desenvolvimento da imaginação, a relação da língua materna com as culturas orais e escritas e a importância da literatura no processo de constituição do \textbf{sujeito}.
Escuta, Fala e Pensamento Aprender a escutar, a organizar o pensamento e a materializá -lo por meio da fala, emitindo ideias, sentimentos, desejos, necessidades, \textbf{posicionamentos} e argumentos, diante de ações, de fatos, de acontecimentos e de conhecimentos é \textbf{uma} atividade essencialmente humana e \textbf{um} dos principais aprendizados da criança.
Escutar, falar e pensar têm origens diferentes no processo de aprendizagens e desenvolvimento da criança, mas \textbf{um} \textbf{influencia} no desenvolvimento do outro e num determinado momento passam a ser coincidentes.
Escuta é o acolhimento e a compreensão do \textbf{que} o outro tem a \textbf{dizer} por meio de suas diferentes expressões – gestual, sonora, oral, escrita e gráfica. \textbf{Exige} sensibilidade, disponibilidade e capacidade de conexão ao outro, por meio da empatia, do desejo e do interesse em interpretar o \textbf{que} está sendo \textbf{dito}.
Compreendida num sentido amplo, envolve o contexto em \textbf{que} a situação comunicativa acontece, a entonação de voz, os olhares, a postura corporal, as vestimentas, o local onde ela ocorre.
Escutar pressupõe, em diferentes situações comunicativas presentes no cotidiano, desenvolver as capacidades de prestar atenção, de antecipar sentidos, de fazer conclusões, de saber respeitar a opinião e a vez do outro, de articular o pensamento e de se expressar para manter o diálogo.
A fala se constitui na transposição de sons, gestos e silêncios, em palavras, em discurso humano ( \textbf{OLIVEIRA}, 2018 ), quer seja na forma oral ou em sinais.
Aprender a usar o vocabulário de \textbf{uma} língua é \textbf{um} processo complexo, porque os gestos, os balbucios devem ser transformados em signos linguísticos, acordados coletivamente, \textbf{que}, organizados em diferentes gêneros textuais, permitem a comunicação entre os \textbf{sujeitos}.
De acordo com Goulart e Mata ( 2016 ), a criança quando chega ao mundo, antes mesmo de nascer, já é referida pelos familiares e pelas pessoas \textbf{que} convivem com a mãe quando conversam com ela ainda na barriga.
Ou seja, desde a concepção, a criança já faz parte de \textbf{um} grupo social, \textbf{que} possui determinados hábitos, costumes e \textbf{uma} língua.
Assim, é nas interações cotidianas com essa língua, \textbf{que} se materializa por meio da palavra, compreendida como meio de comunicação e como \textbf{um} conceito e \textbf{uma} generalização ( VIGOTSKI, 2000 ), \textbf{que} a criança aprende a falar e \textbf{que} o mundo é apresentado, nomeado e significado para e por ela.
Segundo o \textbf{autor} citado, a palavra constitui a consciência e é o elo entre a fala e o pensamento.
O pensamento constitui -se numa ação do intelecto, \textbf{que} possibilita representar mentalmente conhecimentos, fatos, ideias, opiniões, sentimentos, desejos, necessidades, interesses e emoções.
Sua materialização ocorre por meio da palavra, de forma \textbf{que} pensamento e palavra possuem \textbf{uma} relação indissociável e interdependente, mesmo \textbf{que}, em sua origem, os seus processos não sejam coincidentes.
Essa \textbf{inter-relação} entre pensamento e palavra se dá por meio dos significados, \textbf{que} são conceitos estabelecidos socialmente e \textbf{que} permitem a comunicação e o diálogo entre as pessoas.
No entanto, os significados modificam -se ao longo do processo de aprendizagens e desenvolvimento de cada \textbf{sujeito}.
Para \textbf{uma} criança de dois anos o significado de \textbf{uma} palavra é bem diferente do \textbf{que} quando maior ou na fase adulta.
Pode -se citar, por exemplo, a palavra moradia \textbf{que} significa local onde as pessoas habitam para se proteger do sol, do frio ou da chuva.
Para \textbf{uma} criança pequena suas referências serão as casas das pessoas com quem convive, conforme cresce esse conceito é modificado, à medida \textbf{que} estuda sobre os diferentes tipos de moradia da sua época e de outras, conhecendo e visualizando variadas imagens.
Isso ocorre porque os significados são construídos a partir das experiências vivenciadas pelos \textbf{sujeitos} na interação com o outro e com o \textbf{que} significam as palavras em situações concretas.
Como a criança aprende a escutar, a falar e a pensar?
De \textbf{que} forma esses processos são desenvolvidos?
A princípio, a criança se manifesta por meio de choros, balbucios, risadas, gritinhos e expressões faciais.
Essas são suas formas de interagir com o mundo, na busca de apreendê -lo.
Dessa forma, a criança percebe o mundo por meio do olhar, dos movimentos corporais, do toque, do cheiro, da experimentação de diferentes sabores, texturas, \textbf{percepção} dos sons e dos silêncios, do diálogo e da interação com o outro.
Mas, só perceber o mundo não é suficiente, é \textbf{preciso} \textbf{que} alguém o interprete para \textbf{que} ela possa significá -lo.
Transformar o sentir, o perceber, o explorar em linguagem, portanto em significados e sentidos, é o \textbf{que} torna o mundo compreensível para a criança.
Como a criança decifra e interpreta o mundo?
Nas relações sociais.
Desde \textbf{que} nasce ela se empenha em compreender os signos utilizados pelos humanos ao seu redor e procura apropriar -se deles para se fazer entendida pelos outros.
Ao chorar, a mãe, ou o outro adulto \textbf{que} cuida da criança, interpreta esse choro como desconforto \textbf{dizendo} “ Está com fome?
Quer papá, querido? ” ou “ Quer colinho?
Quer sair desse berço?
Venha com a mamãe! ”.
Ao realizar essa ação, transforma em palavras esse incômodo da criança e dá sentido às suas expressões, por meio desse gesto de empatia.
Ou seja, essas expressões ganham intencionalidade a partir das interpretações e reações dos outros perante elas, \textbf{que}, além de traduzir em palavras as sensações do bebê, pode ainda oferecer seu colo, seu carinho, seu olhar, sua raiva e irritação.
Com o tempo e as constantes significações \textbf{que} o adulto vai dando aos gestos, expressões faciais, choros dos bebês no dia a dia, eles vão aprendendo a interpretar suas próprias sensações corporais elaborando representações sobre as situações vivenciadas e desenvolvendo a capacidade de pensar e de se expressar.
Assim, aprender a acalmar, a chorar mais alto, a esperar, a sentir -se querido ou rejeitado, a balbuciar, a produzir as primeiras palavras com significados difusos relativos aos contextos situacionais - qué, dá, bô -, a combinar palavras iniciando o processo de gramaticalização, flexão e sintaxe - qué aua, dá neném, fazi cocô - e, gradativamente, descobrir \textbf{que} tudo tem \textbf{um} nome e \textbf{que} há \textbf{uma} certa estabilidade nos significados das palavras é muito complexo para a criança e, a isso, se denomina o desenvolvimento da fala.
Vygotsky ( 2000 ), ao abordar o desenvolvimento do pensamento e da fala na criança, esclarece \textbf{que} a criança no início do processo de verbalização tem a necessidade de falar \textbf{enquanto} age.
Isto é, \textbf{enquanto} a criança realiza algo, ela \textbf{diz} para si mesma o \textbf{que} está fazendo.
Vygotsky denomina essa fala como organizadora do pensamento e da ação da criança.
À medida \textbf{que} planeja mentalmente suas ações, a criança elabora o seu discurso interior, \textbf{que} se \textbf{configura} como o início do pensamento de forma mais elaborada.
Começa, então, a verbalizar para o outro o \textbf{que} faz.
O \textbf{autor} citado denomina esse processo de fala comunicativa.
Ao distinguir a fala para si da fala para o outro, a criança materializa o pensamento na palavra e, ao mesmo tempo, pode modificar o seu pensamento ao verbalizá -lo.
Assim, apropriar -se da linguagem oral é mais do \textbf{que} apreender \textbf{um} idioma, é apropriar -se de formas de expressão e comunicação, formas estas \textbf{que} constituem o pensamento e a imaginação.
A criança amplia seus conhecimentos, desenvolve novas estruturas linguísticas e aprende a construir narrativas a partir da qualidade, da quantidade e da variedade de situações comunicativas com as quais ela convive.
Nesse processo, passa das nomeações para as descrições de fatos e de situações até \textbf{conseguir} construir narrativas mais elaboradas com \textbf{uma} sequência lógica e autoral.
A criança surda se apropria do significado de expressões, objetos, situações, fatos, acontecimentos, por meio da palavra proferida pelos gestos.
É importante ressaltar \textbf{que}, \textbf{uma} vez detectada a surdez, a criança e a sua família têm o direito de se apropriarem da Libras.
Quanto mais cedo essa língua for inserida no processo de aprendizagens e desenvolvimento da criança, mais possibilidades de significação do mundo ela terá.
A ampliação do vocabulário e o domínio progressivo do uso das palavras em situações comunicativas, utilizando estruturas linguísticas cada vez mais elaboradas, se dará com a promoção de interações com as diferentes linguagens.
Isto é, a criança produz e se apropria da língua quando chamada a falar, a \textbf{dizer} ou contradizer, a escutar/ouvir, a descrever, a contar/narrar, a opinar, a concordar ou discordar, a \textbf{argumentar}, a propor, a definir, a compreender e participar efetivamente do mundo.
Na relação com todas as crianças, sejam surdas, cegas, com paralisia cerebral, transtorno ou síndrome \textbf{que} interfere na produção e apropriação da língua materna, o desafio é construir formas de compreender suas emoções, suas formas de significar e interpretar o mundo e de se fazer compreendido por elas, a fim de garantir o desenvolvimento da linguagem, do pensamento e da imaginação.
Imaginação O \textbf{que} é imaginação?
Qual sua importância para o processo de aprendizagens e desenvolvimento?
Diferente do \textbf{que} se costuma acreditar, a imaginação não é, necessariamente, sinônimo de fantasia, no sentido de pensá -la como algo fora da realidade.
A fantasia se constitui de \textbf{uma} combinação de elementos retirados da realidade \textbf{vivida} \textbf{que} são transformados e reelaborados pela imaginação.
Isso significa \textbf{que} a imaginação é \textbf{um} processo de invenção, de descobertas e de criação \textbf{que}, apoiado na realidade e na produção cultural, ou seja, no \textbf{que} foi \textbf{vivido} anteriormente pelo \textbf{sujeito}, permite a ele fazer recombinações e (re)elaborações produzindo algo novo, quer seja \textbf{um} objeto externo ou \textbf{uma} construção da \textbf{mente} ou da emoção ( VIGOTSKI, 2009 ).
A imaginação é fundamental para a sobrevivência da espécie humana, porque sem ela não seria possível a construção do mundo sociocultural, \textbf{uma} vez \textbf{que} tudo \textbf{que} não pertence ao mundo natural consiste em produção humana desenvolvida a partir da sua atividade criadora.
Por isto, só existem casas, carros, aviões, submarinos, tanques de guerra, armas etc., porque alguém teve a ideia e usou a imaginação para criá -los.
Diante da sua importância, o desenvolvimento da imaginação deve ser \textbf{uma} ação valorizada e garantida nas instituições de Educação Infantil.
Isso se dará à medida \textbf{que} as ações pedagógicas forem as mais variadas possíveis, envolvendo as diferentes linguagens e os conhecimentos do patrimônio da \textbf{humanidade}, assim como a criança ter liberdade para experimentar, manusear diferentes materiais, texturas, (re)organizar espaços e objetos, expressar suas ideias, sendo autorais em seus \textbf{posicionamentos} e em suas produções, quer sejam textos orais, desenhos, escritas etc.
Nessa perspectiva, para a criança imaginar e criar é necessário conhecer, significar, imitar, sentir, cheirar, provar, apalpar, se emocionar.
Desse modo, se alguém propuser a \textbf{uma} criança \textbf{que} imite \textbf{um} cachorro, ela terá \textbf{que} evocar na memória imagens, cheiros, sons emitidos, filmes assistidos e vivências \textbf{que} teve com esse animal, para \textbf{que} possa ter \textbf{uma} imagem mental de \textbf{um} cachorro. \textbf{Logo}, terá \textbf{que} ter conhecido \textbf{um} antes, senão pessoalmente, por meio de representações imagéticas.
Quanto mais experiências tiver com cachorro, mais saberá \textbf{que} é \textbf{um} animal \textbf{que} se apresenta de diferentes formas, pequeno-grande, calmo-agitado, manso-bravo, peludo-pouco pelo.
Assim, a criança terá à sua disposição diferentes elementos para realizar a imitação.
O mesmo ocorre no processo de representação gráfica do cachorro, no desenho, na pintura, em telas, computadores, tablets, papéis etc.
Nas brincadeiras de faz de conta, a imaginação e a fantasia expressam a capacidade da criança de imitar o real, mas não da forma como foi vivenciado, e sim, com vários elementos diferentes, resultantes da sua criação.
Esse processo, de imaginar e de fantasiar, auxilia a criança a compreender a realidade e os diferentes papéis \textbf{que} os \textbf{sujeitos} ocupam nos grupos sociais.
Assim, a criança, ao recordar e reviver suas experiências em diferentes contextos, cria novas combinações e possibilidades de interpretar e representar o real, de acordo com seus desejos, suas necessidades, seus afetos e suas emoções.
Em síntese, a matéria-prima para a imaginação criadora é a memória das experiências \textbf{vividas} proporcionada pela participação efetiva da criança em práticas sociais.
Língua Materna e suas Relações com as Culturas Orais e Escritas A apropriação e a compreensão da língua materna, como já mencionado, é \textbf{um} processo complexo e acontece por meio das interações \textbf{que} ocorrem nas práticas sociais, a partir da utilização da linguagem verbal, \textbf{que} corresponde à linguagem oral e à linguagem escrita.
Essas duas linguagens se \textbf{inter-relacionam}, mas não são iguais, possuem características próprias e saberes específicos \textbf{que} as definem.
Ambas são importantes e necessárias para a construção cultural de \textbf{um} povo e para a organização da vida coletiva. \textbf{Contudo}, \textbf{dependendo} da constituição dos grupos sociais \textbf{uma} linguagem pode ser mais valorizada e utilizada \textbf{que} a outra.
Em determinadas comunidades basta a palavra proferida oralmente para ser realizada \textbf{uma} transação comercial, em outras, é necessário contrato com assinatura para \textbf{que} o negócio se efetive.
Esses diferentes usos da linguagem verbal não podem levar à consideração de \textbf{que} \textbf{uma} linguagem é superior à outra.
É \textbf{preciso} ressaltar esse aspecto, porque, em sociedades em \textbf{que} a escrita é utilizada de forma sistemática há \textbf{uma} tendência em não valorizar a linguagem oral.
Na organização da ação pedagógica das instituições de Educação Infantil, a linguagem verbal deve ser trabalhada com igual importância no processo de aprendizagens e desenvolvimento das crianças, por meio das culturas orais e das culturas escritas.
Mas, o \textbf{que} elas significam?
O \textbf{que} diferencia \textbf{uma} da outra?
Culturas Orais As culturas orais consistem nos saberes, fazeres e conhecimentos produzidos pelos grupos sociais e \textbf{que} são repassados de geração a geração por meio da linguagem oral, compondo sua memória e identidade.
Nos processos de socialização e apropriação desse repertório cultural, são utilizados causos, fábulas, parlendas, provérbios, contos, versos, parábolas, canções, mitos, rezas, representações em folguedos, encenações, como \textbf{uma} forma de entender a realidade, seus fenômenos e acontecimentos.
Esses gêneros textuais, geralmente, são repassados boca a boca e não há \textbf{um} \textbf{autor} em específico, porque à medida \textbf{que} o outro se apropria, elementos são retirados ou acrescentados de acordo com sua realidade, tendo muitas variações para \textbf{um} mesmo texto falado.
Assim, autoria é denominada como de “ domínio \textbf{popular} ”, \textbf{uma} vez \textbf{que} pertence à memória e à cultura de \textbf{um} grupo. \textbf{Uma} das características da linguagem oral é a variedade linguística utilizada pelos \textbf{sujeitos} nas situações comunicativas, devido à diversidade de contextos, de intenções e de usos da língua, na maioria dos casos, a proximidade dos interlocutores pela presença física do outro, possibilitando esclarecer melhor ou repetir o \textbf{que} foi \textbf{dito}, perceber os gestos de quem está falando e a entonação da voz - sugestão, pedido, ordem, ameaça, conversa -, bem como estabelecer \textbf{uma} referência comum para o \textbf{que} está sendo \textbf{dito}.
Diferente da linguagem escrita, em \textbf{que} o texto tem \textbf{que} ser o mais claro possível para \textbf{que} o outro compreenda o significado, porque o leitor não terá a presença física do \textbf{autor}.
A linguagem oral é composta por situações comunicativas informais e também por situações públicas de comunicação oral, em \textbf{que} há registros \textbf{que} a precedem e o monitoramento da fala \textbf{enquanto} é proferida, como em palestras, conferências, mesas redondas, debates, apresentações temáticas, seminários, programas de rádio, de televisão, do YouTube.
Considerar a linguagem oral nessa \textbf{totalidade} significa reafirmar \textbf{que}, tanto em situações comunicativas formais quanto informais, há intencionalidade no \textbf{dizer}, não sendo, portanto, \textbf{uma} ação espontânea.
Na linguagem oral, o escutar, o falar, o repetir, o imitar e o observar são fundamentais, para \textbf{que} os \textbf{sujeitos} sejam capazes de se comunicar.
Nas instituições de Educação Infantil essas aprendizagens podem ser promovidas quando a criança é incentivada a : dialogar com outra ; se expressar nas rodas de conversas ; socializar com outras turmas o \textbf{que} foi aprendido ; transmitir recados ; recontar histórias ; recitar parlendas, versos, poemas, trava-línguas ; brincar de faz de conta, em \textbf{que} ela tem \textbf{que} se colocar diante do outro - posto de gasolina, táxi, pet shop, feiras, caixa de supermercado, médico - ; elaborar, com auxílio do(a) professor, as notícias faladas da turma ou da instituição para serem socializadas com as famílias ; entre tantas outras situações.
A criança \textbf{que} tem a possibilidade de expressar desejos e sentimentos, de desenvolver sua capacidade de argumentação, de dar respostas, de questionar, de avaliar fatos, situações, sem reprimendas dos adultos, desenvolve autoconfiança e maior segurança para se posicionar diante dos outros.
Culturas Escritas Cultura escrita, de acordo com Galvão ( 2016 ), se refere ao lugar simbólico e material \textbf{que} a escrita ocupa em determinado grupo social.
Ou seja, está relacionada à importância e ao valor atribuído pelos \textbf{sujeitos} à escrita, bem como à sua utilização no dia a dia.
A escrita é valorizada de forma diferente pelas pessoas, a \textbf{depender} dos costumes, valores, hábitos e organização dos grupos sociais aos quais elas pertencem, sendo necessário utilizar a terminologia no plural, “ culturas escritas ”, para contemplar essa diversidade.
Por \textbf{que} é importante para o professor compreender o \textbf{que} são as culturas escritas e no \textbf{que} elas \textbf{influenciam} no desenvolvimento da ação pedagógica com as crianças?
A criança, desde bebê, participa de práticas sociais, em \textbf{que} a leitura e a escrita se fazem presentes, porque \textbf{vive} em \textbf{uma} sociedade organizada por e a partir da linguagem escrita, denominada grafocêntrica.
No entanto, como já foi apontado, a forma como essa participação se efetiva está relacionada aos modos como sua comunidade utiliza a escrita.
O \textbf{que} lê e escrevem?
Quando e para \textbf{que} usar a leitura e a escrita?
Quais gêneros textuais são mais comuns e com quais suportes \textbf{opera}?
As respostas a essas questões indicam \textbf{que} em sociedades complexas e heterogêneas, o espaço do escrito envolve relações de poder, é múltiplo, desigual e se transforma permanentemente.
As instituições de Educação Infantil têm \textbf{um} papel fundamental no processo de possibilitar à criança o conhecimento e a participação em diferentes práticas sociais em \textbf{que} o escrito tem \textbf{uma} função real, buscando equiparar suas oportunidades, assim como entender o significado e a importância \textbf{que} o escrito ocupa no dia a dia dela, garantindo a compreensão das culturas escritas, para \textbf{que}, desde pequena se aproprie do como se escreve, para \textbf{que} e para quem se escreve.
Nessa perspectiva, qual é o papel da Educação Infantil no processo de compreensão e apropriação da linguagem escrita?
Nas DCNEI/2009 ( Parecer CEB n. 20 ) é afirmado \textbf{que} a criança tem interesse pela linguagem escrita, antes mesmo do trabalho formal realizado pela instituição de Educação Infantil e apresenta \textbf{que} > Sua apropriação pela criança se faz no reconhecimento, compreensão e fruição da linguagem \textbf{que} se usa para escrever, mediada pela professora e pelo professor, fazendo -se presente em atividades prazerosas de contato com diferentes gêneros escritos, como a leitura diária de livros pelo professor, a possibilidade da criança desde cedo manusear livros e revistas e produzir narrativas e ‘ textos ’, mesmo sem saber ler e escrever. ( p.94 ) Diante dessas indicações, apresentadas num documento de caráter normativo, o papel da Educação Infantil no processo de apropriação e de compreensão da linguagem escrita é de : - propiciar a participação das crianças em práticas sociais de leitura e de escrita de forma a favorecer a compreensão dos seus usos e funções na sociedade ; - proporcionar às crianças, desde bebês, a apreciação e a convivência diária e sistemática com a contação e a leitura de histórias, ampliando seu repertório \textbf{literário} e o conhecimento de \textbf{autores} e ilustradores variados ; - favorecer por meio da organização de tempos, espaços e materiais \textbf{que} as crianças explorem e brinquem, utilizando objetos \textbf{que} fazem parte das culturas escritas, como embalagens, cadernetas, livros, \textbf{agendas}, lápis, canetas, revistas, jornais, folders, caixas, bulas de remédios, cadernos ; - possibilitar às crianças, mediadas pelo professor, a leitura e a escrita de forma espontânea[^20 ] e convencional, de gêneros textuais variados, em \textbf{que} elas possam compreender sua estrutura, bem como os diferentes suportes \textbf{que} podem ser veiculados, inclusive o digital, de acordo com a intenção comunicativa ; - aguçar o interesse e a curiosidade das crianças pela leitura e pela escrita, por meio de situações comunicativas reais, em \textbf{que} elas possam iniciar a compreensão das regras do Sistema de Escrita Alfabético, como por exemplo, \textbf{que} se escreve de cima para baixo, da esquerda para a direita, \textbf{que} são utilizadas letras para escrever e não outros símbolos etc.
Portanto, é necessário reafirmar \textbf{que} a Educação Infantil tem o papel de oportunizar à criança a compreensão dos diversos usos e funções \textbf{que} a escrita tem na sociedade, a \textbf{percepção} de regularidades, como as citadas e a estrutura formada por letras e sinais gráficos.
Cabe ainda ressaltar \textbf{que} a linguagem escrita deve ser trabalhada com a mesma importância das demais linguagens – visuais, corporal, gráfica, musical, audiovisual etc. – em contextos reais e significativos.
A \textbf{tarefa} do Ensino Fundamental é a alfabetização stricto sensu, sendo compreendida como processo de apropriação, compreensão e efetivo uso do complexo Sistema de Escrita Alfabético a partir de \textbf{um} trabalho pedagógico \textbf{explícito} e sistemático no entendimento do seu funcionamento, suas regularidades e irregularidades.
Diante da definição do papel da Educação Infantil, cabe perguntar : como as crianças participam de práticas de leitura e de escrita?
Crianças pequenas lêem e escrevem, o \textbf{que} e como lêem e escrevem?
Todas as crianças lêem o mundo desde o ventre da mãe.
Lêem vozes, sons, expressões faciais, cheiros, temperatura do ambiente, humor de quem cuida delas, colo, gestos de carinho, afago, irritação e rispidez.
À medida \textbf{que} essas leituras vão sendo feitas, elas buscam interpretá -las e vão atribuindo sentidos ao \textbf{que} lhes acontecem.
Quando, no momento da troca, o professor se coloca à disposição da criança, mantém contato visual, acolhe os gestos \textbf{que} ela faz de estender o braço e tentar tocar o seu rosto, e esse se \textbf{aproxima}, confirmando o toque \textbf{que} ela busca e conversa com a criança, ele estabelece \textbf{um} diálogo e possibilita a descoberta da linguagem.
Todas as expressões, visuais, táteis ou sonoras, \textbf{que} forem ofertadas às crianças funcionam como ‘ livro ’, material a ser lido e interpretado pela criança e pelo adulto.
Assim, compreende -se leitura num sentido amplo, em \textbf{que} a leitura de mundo, como já afirmava Paulo Freire, precede a leitura da palavra, sendo a leitura relacionada à compreensão dos signos linguísticos, extensão e consequência, dessas diferentes leituras.
Nesse sentido, é importante ressaltar \textbf{que} outros gêneros textuais e suportes têm surgido, com o uso cada vez mais frequente das tecnologias digitais no dia a dia, \textbf{exigindo} novas formas de leitura.
É função das instituições de Educação Infantil e dos ( as ) professores ( as ) proporcionarem à criança, desde bebê, \textbf{que} participe de variadas práticas de leituras e de registros, em \textbf{que} tenha possibilidade de expressar opiniões sobre o \textbf{que} gostou ou não, qual personagem gostaria de ser, se fosse o \textbf{autor} qual final seria dado à história, de brincar com palavras por meio de rimas e aliterações, de conversar sobre quais estratégias utiliza para ganhar em jogos eletrônicos, entre outras.
De acordo com Vygotsky ( 2000 ), a escrita por ser \textbf{uma} representação simbólica não pode ser compreendida de forma dissociada de outras linguagens.
O seu processo de apropriação tem origem muito antes da criança \textbf{conseguir} escrever de forma convencional utilizando letras, porque a produção de textos por meio de signos inicia -se pela criança a partir da utilização do próprio corpo com gestos indicativos e expressões corporais \textbf{que} são interpretadas pelo outro e \textbf{que} posteriormente, são incorporados aos desenhos e às brincadeiras de faz de conta.
Essas diferentes representações utilizadas pela criança ( gestos, desenho, faz de conta ), somadas às vivências com pinturas, colagens, modelagens, utilização de letras e outros símbolos, imagens e situações \textbf{vividas}, observadas ou imaginadas lhes possibilitam criar textos orais \textbf{que} podem ser escritos por ela de forma espontânea ou ditados para \textbf{um} escriba transcrever.
De acordo com essa perspectiva, a criança, desde bebê, é capaz de produzir textos, de deixar \textbf{marcas}, portanto, lê e escreve de formas variadas, até o momento em \textbf{que} suas \textbf{marcas} se constituem em signos linguísticos acordados socialmente expressos na forma de \textbf{uma} escrita convencional.
Para tanto, é \textbf{preciso} \textbf{que} a criança perceba a importância e a necessidade de aprender a escrever, sendo fundamental \textbf{que} ela participe de situações comunicativas em \textbf{que} a linguagem escrita é utilizada, como : escrita de bilhetes para as famílias ; identificação de brinquedos, de objetos da sala e de pertences pessoais ; produção de cartazes para socializar com outras crianças o \textbf{que} está sendo aprendido ; lista de histórias favoritas ; produção de poemas para recitar num sarau ; produção de livro da turma com recontos de histórias ; entre outras.
Literatura A literatura se apresenta como possibilidade do leitor vivenciar a alteridade, deslocar -se para \textbf{viver} outras vidas, outros mundos imaginários novos e surpreendentes.
Possibilita, assim, \textbf{um} conhecimento mais profundo do humano ao \textbf{aproximar} o leitor dos sentimentos, das emoções, das situações, dos questionamentos, das inquietações trazidas pelo outro ( personagem, contexto ), enriquecendo sua própria experiência.
Pela capacidade de potencializar processos de humanização, a literatura será tratada nesse documento como direito inalienável do \textbf{sujeito} ao acesso, fruição e apropriação.
Ítalo Calvino ( 1990 ) aponta algumas características da literatura \textbf{que} auxilia o leitor, nesse caso a criança, no processo de humanização : - a leveza do texto \textbf{literário} é a possibilidade de tratar de temas complexos e difíceis como a violência, a morte, a separação, o abandono, o desprezo, de forma artística, colocando a criança na dimensão do sensível, permitindo \textbf{que} ela elabore o assunto a partir de \textbf{uma} outra lógica ; - a exatidão é a capacidade \textbf{que} o texto \textbf{literário} tem de organizar as palavras de forma concisa e \textbf{precisa} para expressar pensamentos, \textbf{raciocínios}, ideias, emoções ; - a rapidez está vinculada ao tempo \textbf{literário}, \textbf{que} não é cronológico, porque pode se passar muitos anos de \textbf{um} parágrafo para outro ou pode imputar à criança a espera para saber o \textbf{que} vai acontecer.
Apresenta, também, sucessão, ordem, simultaneidade ; - a visibilidade é a característica \textbf{que} oportuniza a imaginação do narrado, a criação de cenários e a corporificação de personagens, \textbf{despertando} -lhe sentimentos e sensações reais.
Também inclui a relação entre o texto e a ilustração \textbf{que} provoca leituras dos não ditos e novas interpretações ; - a multiplicidade \textbf{revela} os mistérios presentes na literatura como ser \textbf{um} e muitos ao mesmo tempo, ser a princesa ou o príncipe e ser a bruxa malvada. \textbf{Viver} várias vidas se tornando outros, enriquecendo a própria existência ; - a consistência está relacionada à qualidade do \textbf{que} é oferecido ao leitor, as escolhas dos temas, dos livros e das formas como serão tratados.
Nesse sentido, pensar o trabalho com a literatura nas instituições de Educação Infantil \textbf{implica} superar \textbf{uma} visão instrucional e pragmática, na qual lê -se para ensinar as letras ou para a criança aprender a se comportar.
Compreendê -la como lugar : de relações ; de brincadeiras ; de produção de sentido e significados compartilhados ; de conhecimento de si e do outro ; de constituição da subjetividade e de sentimento de pertença ; de ampliação das experiências com a cultura escrita ; é a possibilidade de efetivar \textbf{um} trabalho diferenciado com a literatura, na perspectiva de garantir esse direito inalienável do \textbf{sujeito}, de se tornar mais humano por meio de escritos e imagens de outrem.
A literatura é capaz de deslocar o leitor/ouvinte quando provoca emoções, risos e choros, quando promove o diálogo entre o leitor/ouvinte e a obra por meio das ilustrações, da forma como o texto é organizado no livro, das expressões, metáforas, rimas e aliterações usadas pelo \textbf{autor}, provocando efeitos de sentidos.
Em síntese, a leitura \textbf{literária} não é adorno, nem entretenimento para acalmar e divertir as crianças, muito menos \textbf{um} servir de pretexto ou instrumento de ensino da língua escrita ou de outros conteúdos.
Ela possibilita o conhecimento de si e do outro, a ampliação do \textbf{que} já sabe do/sobre o mundo, a organização de experiências pela palavra, a imaginação, a criação, o pensamento, a emoção, a apreciação estética dos textos verbais e visuais e suas interlocuções. [ ^19 ] : O uso dos termos “ fala ” ou “ palavra ” neste documento se refere à expressão do \textbf{sujeito} por meio da linguagem oral ou da Língua Brasileira de Sinais. [ ^20 ] : A escrita espontânea se refere aos registros gráficos das crianças – rabiscos, garatujas, bolinhas, símbolo, números, letras – realizados de forma mais livre e com a intenção de escrever.
Nesse processo, elas mobilizam seus conhecimentos para compreender como se escreve, elaborando ideias, \textbf{teorias} a respeito do funcionamento da escrita alfabética.

% matched lemmas: agenda, aproximar, argumentar, autor, central, configurar, conseguir, contudo, depender, despertar, dizer, enquanto, exigir, explícito, humanidade, implicar, inerente, influenciar, inter-relacionar, inter-relação, literário, logo, marca, mente, necessitar, oliveira, operar, percepção, popular, posicionamento, preciso, que, raciocínio, revelar, sujeito, tarefa, teoria, totalidade, um, viver
\end{textsample}
